\par\noindent\textbf{(1)} 每个\(\delta\)点波函数连续而导数跃变 \(\psi'(0^+)-\psi'(0^-)=(2mg/\hbar^2)\psi(0)\)。用状态 \((\psi,\psi')^T\) 表示，自由段传播矩阵 \(P(a)=\begin{pmatrix}\cos ka&\sin ka/k\\-k\sin ka&\cos ka\end{pmatrix}\)，\(\delta\)点跳跃矩阵 \(D=\begin{pmatrix}1&0\\2mg/\hbar^2&1\end{pmatrix}\)。
\par\noindent\textbf{(2)} 单胞矩阵为 \(M=DP(a)\)，两矩阵行列式均为1。Bloch本征值 \(e^{\pm iqa}\) 给 \(\operatorname{tr}M=2\cos qa\)，故 \(\cos(qa)=\cos(ka)+(mg/(\hbar^2k))\sin(ka)\)。
\par\noindent\textbf{(3)} 定义 \(F(k)=\cos(ka)+(mg/(\hbar^2k))\sin(ka)\)。当 \(|F(k)|\le1\)，存在实准动量 \(q\) 传播Bloch态；\(|F(k)|>1\) 时 \(q\) 复数，形成禁带。该禁带不是由周期\(\delta\)势中的吸收产生，模型Hamilton量仍为自伴且无耗散。
\par\noindent\textbf{(4)} 当 \(g\to0\)，\(F(k)\to\cos ka\)，全部正能连续自由粒子色散只发生布里渊区折叠，禁带闭合。小 \(g>0\) 时在Bragg点 \(ka\approx n\pi\) 附近打开小禁带，反映向前和反向平面波的简并耦合。
